\section{Grundlagen und Stand der Technik}\label{sec:grundlagen}

\subsection{Jass und Jasskarten}\label{sec:jass}

Ein Jassblatt hat 36 Karten in vier Farben mit je neun Werten.
Diese Arbeit verwendet zwei Blätter von AGMüller: ein französisches
Standardblatt und das deutschschweizerische \enquote{Schaffhauser
Spielkarten Opti-Quatro}. Eine Partie wird mit einem Blatt gespielt; die
beiden werden nicht gemischt. Viele Karten sehen nach einer Drehung um
180\textdegree{} fast gleich aus. Spiegeln verändert dagegen die Anordnung
der Eckzeichen und eignet sich deshalb nicht als Augmentierung.

Wie viele Punkte eine Karte zählt, hängt von der Spielart und von der
Trumpffarbe ab. Bei den Spielarten, welche die App unterstützt, ergeben
alle Karten zusammen 152 Punkte. Für den letzten Stich kommen weitere
5 Punkte hinzu \cite{swisslos2026grundlagen}. Bei anderen Spielarten
können auch bestimmte Stiche oder Kartenkombinationen zusätzliche Punkte
geben. In dieser Arbeit stehen die Erkennung der obersten Karte und die
Berechnung der Punkte eines Stapels im Mittelpunkt.

\subsection{Objekterkennung mit Deep Learning}\label{sec:objekterkennung}

Bildklassifikation ordnet einem Bild eine Klasse zu. Objektdetektion gibt
zusätzlich die Position als Bounding Box an. Einstufige Detektoren wie YOLO
sagen Box und Klasse in einem Durchgang voraus \cite{redmon2016yolo}. Ein
zweistufiger Ansatz lokalisiert zuerst und klassifiziert danach. In dieser
Arbeit werden beide Ansätze und die reine Klassifikation verglichen
(Abschnitt~\ref{sec:varianten}).

\subsection{Die YOLO-Familie und YOLO11}\label{sec:yolo}

Verwendet wird YOLO11 mit Ultralytics \cite{jocher2024ultralytics,
ultralytics2025yolo11docs,khanam2024yolov11}. Für die drei Ansätze werden verschiedene
Ausführungen derselben Modellfamilie trainiert: A und B$_2$ klassifizieren,
B$_1$ lokalisiert mit einer gedrehten Box, und C liefert eine
achsenparallele Box samt Kartenklasse. Alle nutzen die kleinste
Modellgrösse n, damit sie auf einem Smartphone laufen. Die
Mosaic-Augmentierung setzt beim Training vier Bilder zusammen und wird
in den letzten Epochen abgeschaltet \cite{bochkovskiy2020yolov4,ultralytics2025augment}.

YOLO26 erschien erst im Januar 2026 und war damit zu Beginn der
Umsetzung noch sehr neu. Das Modell kann Vorhersagen ohne
nachgelagerte Non-Maximum Suppression ausgeben. Dadurch unterscheidet
sich sein Ausgabeformat von YOLO11 \cite{ultralytics2026yolo26}. Die
Unterstützung dieses Formats war in mobilen Anwendungen zu diesem
Zeitpunkt noch nicht vollständig ausgereift. Ein GitHub-Issue vom
Mai 2026 zeigt beispielsweise, dass die Verarbeitung dieser Ausgabe
noch ergänzt werden musste \cite{ultralytics2026flutter487}.

Für diese Arbeit war eine verlässliche Integration auf iOS und Android
wichtiger als der Einsatz der neuesten Modellversion. YOLO11 liess sich
bereits stabil trainieren, exportieren und in beiden Apps ausführen.
Deshalb wurde das Projekt mit YOLO11 ausgeführt.

\subsection{Metriken}\label{sec:metriken}

Die Modelle werden mit verschiedenen Kennzahlen verglichen. Bei einem
Detektor zählt neben der richtigen Kartenklasse auch die Position. Die
\emph{Intersection over Union} (IoU) misst, wie stark sich die
vorhergesagte Box und die von Hand gesetzte Box überlappen: 0 bedeutet
keine Überlappung, 1 eine deckungsgleiche Box.

Die \emph{Precision} gibt an, wie viele der gemeldeten Karten richtig
sind. Meldet ein Modell zehn Karten und acht davon sind richtig, beträgt
die Precision 80\,\%. Der \emph{Recall} beschreibt, wie viele der
vorhandenen Karten gefunden werden. Sind zehn Karten vorhanden und das
Modell findet acht davon, beträgt der Recall ebenfalls 80\,\%. Beide
Werte hängen von der gewählten Konfidenzschwelle ab.

Die \emph{Average Precision} (AP) fasst Precision und Recall für eine
Kartenklasse zusammen. Die \emph{mean Average Precision} (mAP) bildet
daraus den Mittelwert über alle Klassen. Bei \map{50} zählt eine
Vorhersage nur dann als richtig, wenn sie mindestens eine IoU von 0.5
erreicht \cite{everingham2010pascal}. \map{50--95} verwendet mehrere
IoU-Schwellen zwischen 0.5 und 0.95. Ungenaue Boxen werden dadurch
strenger bewertet
\cite{lin2014coco,coco2015eval,ultralytics2025metrics}. Die mAP ist
deshalb nicht mit dem Anteil korrekt benannter oberster Karten
gleichzusetzen.

Für die reine Bildklassifikation wird die \emph{Top-1-Genauigkeit}
verwendet: Der Klassenname mit dem höchsten Wert muss stimmen. Für die
App ist zusätzlich die einfachere Zahl \enquote{oberste Karte richtig}
wichtig. Sie gibt an, auf wie vielen realen Fotos das Modell die oberste
Karte korrekt benennt. Ob die App sie tatsächlich zählt, hängt
anschliessend noch von
Konfidenzschwelle und Stabilitätsregel ab
(Abschnitt~\ref{sec:betriebspunkt}).

\subsection{Transfer Learning}\label{sec:transfer}

Alle Ansätze starten mit vortrainierten YOLO11n-Gewichten, statt ein Netz
von Grund auf zu trainieren \cite{pan2010transfer}. Die Startgewichte
von C stammen aus dem Training auf COCO-Alltagsbildern \cite{lin2014coco}, jene von B$_1$
aus DOTA-Luftbildern \cite{xia2018dota} und jene der Klassifikatoren A und B$_2$ aus
ImageNet \cite{russakovsky2015imagenet,ultralytics2025yolo11docs}. So können bereits gelernte
Bildmerkmale übernommen werden. Die Ausgabe wird für die
projektspezifischen Kartenklassen neu trainiert; deren Anzahl ist je
Ansatz verschieden (Abschnitt~\ref{sec:varianten}).

\subsection{Synthetische Trainingsdaten und Domain Randomization}\label{sec:synthdata}

Gerenderte Bilder haben den Vorteil, dass ihre Labels automatisch und
genau erzeugt werden können. Die Herausforderung liegt bei der
Übertragung auf reale Fotos. Dieser Unterschied wird als
Sim-to-Real-Gap bezeichnet.

Domain Randomization soll diese Lücke verkleinern. Dabei werden
Perspektive, Licht, Textur und Kameraeinstellungen bewusst verändert.
Das Modell sieht dadurch viele unterschiedliche Situationen und soll
lernen, welche Merkmale auch auf realen Bildern wichtig sind
\cite{tobin2017domain}. Dafür muss nicht jedes gerenderte Bild möglichst
echt aussehen. Auch einfache synthetische Bilder können für das Training
geeignet sein, wenn die Objekte plausibel in die Szene eingefügt werden
\cite{dwibedi2017cut}.

Tremblay et al. trainierten einen Detektor nur mit synthetischen Bildern
und konnten damit Autos auf realen Aufnahmen erkennen. Wenige zusätzliche
reale Trainingsbilder verbesserten das Ergebnis weiter
\cite{tremblay2018training}. Ein ähnlicher Ansatz wurde bereits für
Spielkarten verwendet. Das Projekt von geaxgx erzeugt aus fotografierten
Pokerkarten einen synthetischen Datensatz und trainiert damit YOLOv3.
Erkannt werden dort die Eckzeichen der Karten. Eine Auswertung auf realen
Fotos wird jedoch nicht ausgewiesen \cite{geaxgx2018cards}.

Ultralytics empfiehlt mindestens 1500 Bilder je Klasse
\cite{ultralytics2025tips}. Bereits für ein Blatt mit 36 Karten ist diese
Menge von Hand kaum zu erstellen und zu labeln. In dieser Arbeit werden
die Trainingsbilder deshalb gerendert. Die Prüfung erfolgt auf eigenen
realen Aufnahmen (Abschnitt~\ref{sec:validierung}).

\subsection{On-Device-Inferenz}\label{sec:ondevice}

Core ML und Vision führen das exportierte Modell auf iOS aus
\cite{apple2025coreml,apple2025vision}. Auf Android übernimmt LiteRT diese
Aufgabe \cite{google2026litert}. Die beiden Exporte stammen aus denselben Gewichten. Die
Laufzeiten unterscheiden sich bei der Verarbeitung der Boxen und
Konfidenzen, was für den Vergleich der App-Ausgaben geprüft werden muss
(Abschnitt~\ref{sec:betriebspunkt}).
